\newpage
\phantomsection
\label{prf:rational-quotient-zero-iff-numerator-zero}

\begin{remark*}[Return]
\hyperref[thm:rational-quotient-zero-iff-numerator-zero]{Return to Theorem}
\end{remark*}

\begin{theorem*}[A Rational Quotient Is Zero Exactly When Its Numerator Is Zero]
For every $a,b\in\mathbb{Q}$ with $b\neq 0$,
\[
\frac{a}{b}=0
\quad\Longleftrightarrow\quad
a=0.
\]
\end{theorem*}

\begin{proof}
\textbf{Professional Standard Proof.}~\\
Let $a,b\in\mathbb{Q}$ with $b\neq 0$.

First suppose that
\[
\frac{a}{b}=0.
\]
By the definition of rational division,
\[
\frac{a}{b}=a\cdot b^{-1}.
\]
Hence
\[
a\cdot b^{-1}=0.
\]
Since $b\neq 0$, its multiplicative inverse $b^{-1}$ exists and is nonzero.
Because $\mathbb{Q}$ is an integral domain, the zero-product property gives
\[
a=0
\qquad\text{or}\qquad
b^{-1}=0.
\]
Since $b^{-1}\neq 0$, it follows that
\[
a=0.
\]

Conversely, suppose that
\[
a=0.
\]
Then, by the definition of rational division,
\[
\frac{a}{b}=a\cdot b^{-1}=0\cdot b^{-1}=0.
\]
Therefore
\[
\frac{a}{b}=0
\quad\Longleftrightarrow\quad
a=0.
\]
\end{proof}

\begin{proof}
\textbf{Detailed Learning Proof.}~\\

\textbf{Step 1.} Rewrite the quotient as multiplication by an inverse.

Let $a,b\in\mathbb{Q}$ with
\[
b\neq 0.
\]
By the definition of rational division,
\[
\frac{a}{b}=a\cdot b^{-1}.
\]
Thus the statement
\[
\frac{a}{b}=0
\]
is the same as the statement
\[
a\cdot b^{-1}=0.
\]

\textbf{Step 2.} Prove the forward implication.

Assume
\[
\frac{a}{b}=0.
\]
Then
\[
a\cdot b^{-1}=0.
\]
Since $b\neq 0$, the inverse $b^{-1}$ exists. Also,
\[
b^{-1}\neq 0,
\]
because if $b^{-1}=0$, then multiplying by $b$ gives
\[
1=b\cdot b^{-1}=b\cdot 0=0,
\]
contradicting the field axioms.

Since $\mathbb{Q}$ is an integral domain, a product in $\mathbb{Q}$ is
zero only if one of its factors is zero. Therefore, from
\[
a\cdot b^{-1}=0,
\]
it follows that
\[
a=0
\qquad\text{or}\qquad
b^{-1}=0.
\]
The second possibility is impossible, so
\[
a=0.
\]

\textbf{Step 3.} Prove the reverse implication.

Assume
\[
a=0.
\]
Then, by the definition of rational division,
\[
\frac{a}{b}=a\cdot b^{-1}.
\]
Substituting $a=0$ gives
\[
\frac{a}{b}=0\cdot b^{-1}.
\]
By the absorbing property of zero,
\[
0\cdot b^{-1}=0.
\]
Hence
\[
\frac{a}{b}=0.
\]

\textbf{Step 4.} Conclude the biconditional.

The forward implication proves
\[
\frac{a}{b}=0 \;\Longrightarrow\; a=0.
\]
The reverse implication proves
\[
a=0 \;\Longrightarrow\; \frac{a}{b}=0.
\]
Therefore
\[
\frac{a}{b}=0
\quad\Longleftrightarrow\quad
a=0.
\]
\end{proof}

\begin{remark*}[Proof structure]
The proof is a direct proof of both implications. The forward direction
rewrites the quotient as $a\cdot b^{-1}$ and uses the zero-product
property in the integral domain $\mathbb{Q}$. The reverse direction
substitutes $a=0$ into the definition of division and uses the absorbing
property of zero.
\end{remark*}

\begin{remark*}[Dependencies]~\\
\begin{itemize}
  \item \hyperref[def:rational-division]{Definition~\ref*{def:rational-division}}
  \item \hyperref[thm:q-is-integral-domain]{Theorem~\ref*{thm:q-is-integral-domain}}
  \item \hyperref[thm:q-is-field]{Theorem~\ref*{thm:q-is-field}}
  \item Absorbing property of zero under multiplication
\end{itemize}
\end{remark*}
